\section{The Dirac Equation, the Light Cone, and Mass}

A free tone streams one dock per beat and never spreads, so the substrate has a built-in speed limit, one site-step per tick, the same in every direction. That bare fact, sharpened by the way the knit mixes the spinor components as a tone moves, is enough to reproduce relativistic matter. Taken as a quantum walk, the directional propagation of the tone has, in its long-wavelength limit, exactly the Dirac dispersion, the relation between energy and momentum that governs a relativistic spin-half particle. The light cone that comes with it is finite and frame-independent, with a dynamical exponent $z = 1$, meaning a disturbance spreads ballistically and linearly rather than diffusively. This is the Dirac quantum walk realized on the actual mesh, a known construction reproduced on the substrate rather than discovered here (depth L2).

A discrete lattice usually betrays itself by a preferred direction, a faint cubic grain in how fast signals travel along the axes versus the diagonals. The mesh does not, and the reason is triality. Because the twenty-four directions carry the threefold symmetry of $F_4$, the leading anisotropy that a generic lattice would show is cancelled, and the propagation is isotropic to fourth order, measured directly on the substrate. Boosts act as they should, and the cone tilts without changing shape. The flat cusp behaves, to the precision reachable in simulation, like the smooth Lorentzian space of textbook relativity, with the discreteness hidden far below in the ultraviolet.

\subsection{The g-factor, measured not assumed}

A relativistic electron has a magnetic moment, and the Dirac theory famously predicts the dimensionless number that sets it, the $g$-factor, to be exactly $2$. In the standard theory this comes from the structure of the Dirac equation. Here it is not put in, it is read out. Coupling the substrate fermion to a magnetic field produces a Landau spectrum, a ladder of energy levels, and the value of $g$ is fixed by whether that ladder has a state pinned at zero energy. The substrate spectrum has the zero mode, and the number it implies is $g = 2$. The discriminating control is a scalar excitation built on the same substrate, which lacks the zero mode and would have given a different value, so the result could have failed and did not. The $g$-factor of $2$ is therefore a measurement on the crystal, not an assumption (depth L2).

\subsection{Where mass comes from}

A massless fermion is a pure chirality, a left-handed or right-handed object travelling at the speed limit. Mass is what couples the two handednesses together so that the particle can sit still, and in the substrate the two handednesses are nothing other than the two spinor octets, $8s$ and $8c$. Mass is the strength of the coupling between them. When that coupling is switched off the octets decouple, the two chiralities separate, and the fermion becomes a massless Weyl particle. When it is on, the fermion acquires a rest frame and a mass.

This mechanism is measured on the substrate in two independent ways that agree, once from the gap that opens in the dispersion relation and once from the commutator of the chirality operator with the dynamics, and both vanish together in the Weyl limit when the octets are uncoupled (depth L2). What the substrate gives, then, is the \emph{mechanism} of mass, exactly and measurably, namely the coupling of the two spinor octets that triality relates. What it does not yet give is the \emph{value}, the actual numbers in the mass spectrum, which sit at the same honest frontier as the strength of the forces and are taken up with the open problems. The how is in hand. The how-much is not.
